\scanpage{181}
\mspage{16}
% archive sheet number: 180
% block: top, line 1 | text
générateur topologique de $\mathbb Z'(1)$) \deleted{\unsure{pour les}}\inserted{aux} \unsure{différents} $x_i$.

% block: upper third, lines 2-6 | text
\deleted{\unsure{La structure de}} Comme on \unsure{connait|connaît} la structure de $\pi_1^{\unsure{(p)'|p'}}(X_{\overline K},\unsure{\zeta|\xi})$ \inserted{en terme du genre de $X_{\overline K}$ et du nbre de pts \unsure{enlevés}} (SGA 1 XIII) \deleted{\unsure{cela se lit}} on voit que le pb de déterminer la forme des $h_i$ devient assez concret. Cela \deleted{\unsure{pourra être à cet}} sera fait par Deligne \uline{par voie transcendante} dans un exposé ultérieur.

% block: upper to middle third, lines 7-10 | text
Lorsque $S$ est lui-\unsure{$=$} un trait, donc \deleted{$\unsure{\pi_1^{\prime}}$} $I_t=\pi_1(\eta,\overline\eta)_t\simeq\mathbb Z'_{\deleted{\unsure{\ell}}}(1)$, alors l'action extérieure de $I_t$ sur $\pi_1^{(p)'}(X_{\overline K},\zeta)$ est connue quand on \unsure{connait|connaît} l'hom induit $\deleted{\mathbb Z'(1)\mathbin{-}}\pi_1(\eta,\overline\eta)_{\mathrm{tam}}\longrightarrow\pi_1(U_0,\overline\eta)$ i.e. un \unsure{l'hom}

% block: middle third, lines 11-12 | display
\[
\begin{aligned}
&\deleted{\mathbb Z'(1)\mathbin{-}}\\
u:\ &\mathbb Z'(1)\longrightarrow\mathbb Z'(1)^I
\end{aligned}
\]

% block: middle third, lines 13-15 | text
\deleted{dont la connaissance \unsure{équivaut} d'ailleurs, comme on \unsure{voit aussitôt, à la}} qui est d'ailleurs donné par

% block: middle third, line 16 | display
\[
u(\lambda)=\deleted{\unsure{\sum}}(n_i\lambda)\deleted{\unsure{\text{avec}}}
\]

% block: middle to lower third, line 17 | display
\[
\text{avec}\qquad n_i=\operatorname{long}(V/\unsure{\mathfrak a_i|\mathfrak m_i}V).
\]

% block: lower third, lines 18-22 | text
où $\mathfrak a_i$ est l'idéal définissant le diviseur $D_i$\inserted{de $S$}. (L'entier $n_i$ mesure, \deleted{\unsure{en un sens}} heuristiquement, la vitesse avec laquelle \unsure{disparait|disparaît} la singularité $x_i$ en passant de la fibre spéciale à la fibre \unsure{générale|générique}.) Donc le générateur choisi de $\mathbb Z'(1)$ agit \inserted{extérieurement} sur $\pi_1^{p'}(X_{\overline K};\zeta)$ par

% block: bottom, line 23 | display
\[
h=h_1^{n_1}\cdots h_{\unsure{r|\rho}}^{n_{\unsure{r|\rho}}}
\]
